\documentclass[11pt, a4paper]{report}

% ========== Common Packages ==========
\usepackage{fontspec}
\usepackage{kotex}
\usepackage{amsmath, amssymb, amsthm}
\usepackage{mathtools}
\usepackage{bm}
\usepackage{graphicx}
\usepackage{booktabs}
\usepackage{array}
\usepackage{tabularx}
\usepackage{longtable}
\usepackage{hyperref}
\usepackage{cleveref}
\usepackage{geometry}
\usepackage{fancyhdr}
\usepackage{titlesec}
\usepackage{enumitem}
\usepackage{xcolor}
\usepackage{tcolorbox}
\usepackage{algorithm}
\usepackage{algorithmic}
\usepackage{tikz}
\usetikzlibrary{arrows.meta, positioning, calc, decorations.pathreplacing, shapes.geometric, fit, patterns, angles, quotes, trees}
\usepackage{pgfplots}
\pgfplotsset{compat=1.18}
\usepackage{caption}
\usepackage{subcaption}
\usepackage{float}

\geometry{margin=2.5cm, top=3cm, bottom=3cm}

\usepackage{multirow}
\usepackage{siunitx}
% ========== Colors ==========
% Common
\definecolor{defblue}{RGB}{30, 80, 150}
\definecolor{thmgreen}{RGB}{20, 110, 60}
\definecolor{noteorange}{RGB}{200, 100, 0}
\definecolor{lightblue}{RGB}{230, 240, 255}
\definecolor{lightgreen}{RGB}{230, 255, 235}
\definecolor{lightorange}{RGB}{255, 245, 230}
\definecolor{lightgray}{RGB}{245, 245, 245}
% Extended
\definecolor{lightred}{RGB}{255, 235, 235}
\definecolor{darkred}{RGB}{180, 30, 30}
\definecolor{biogreen}{RGB}{10, 130, 80}
\definecolor{cvpurple}{RGB}{100, 40, 150}
\definecolor{injuryred}{RGB}{200, 40, 40}
\definecolor{codegray}{RGB}{240, 240, 245}
\definecolor{pipeblue}{RGB}{25, 100, 180}
\definecolor{constraintred}{RGB}{200, 50, 50}
\definecolor{termblack}{RGB}{30, 30, 30}
\definecolor{goldcolor}{RGB}{180, 140, 20}

% ========== Theorem-like environments ==========
\theoremstyle{definition}
\newtheorem{definition}{Definition}[chapter]
\newtheorem{example}{Example}[chapter]

\theoremstyle{plain}
\newtheorem{theorem}{Theorem}[chapter]
\newtheorem{proposition}{Proposition}[chapter]

\theoremstyle{remark}
\newtheorem{remark}{Remark}[chapter]

% ========== tcolorbox environments ==========
% --- Common (all parts) ---
\newtcolorbox{keyidea}[1][]{
  colback=lightblue,
  colframe=defblue,
  fonttitle=\bfseries,
  title={Key Idea},
  #1
}

\newtcolorbox{importantnote}[1][]{
  colback=lightorange,
  colframe=noteorange,
  fonttitle=\bfseries,
  title={Important},
  #1
}

\newtcolorbox{summary}[1][]{
  colback=lightgreen,
  colframe=thmgreen,
  fonttitle=\bfseries,
  title={Summary},
  #1
}

% --- Warning/Error ---
\newtcolorbox{warningbox}[1][]{
  colback=lightred,
  colframe=darkred,
  fonttitle=\bfseries,
  title={Warning},
  #1
}

% --- Domain: Biomechanics & Clinical ---
\newtcolorbox{clinicalnote}[1][]{
  colback=lightred,
  colframe=darkred,
  fonttitle=\bfseries,
  title={Clinical Note},
  #1
}

\newtcolorbox{biomechbox}[1][]{
  colback=lightgreen,
  colframe=biogreen,
  fonttitle=\bfseries,
  title={Biomechanics},
  #1
}

\newtcolorbox{injurybox}[1][]{
  colback={injuryred!8},
  colframe=injuryred,
  fonttitle=\bfseries,
  title={Injury Risk},
  #1
}

% --- Domain: Computer Vision ---
\newtcolorbox{cvbox}[1][]{
  colback={cvpurple!5},
  colframe=cvpurple,
  fonttitle=\bfseries,
  title={Computer Vision},
  #1
}

% --- Pipeline & Setup ---
\newtcolorbox{setupbox}[1][]{
  colback=lightgray,
  colframe=gray!60,
  fonttitle=\bfseries,
  title={Setup},
  #1
}

\newtcolorbox{pipelinebox}[1][]{
  colback=pipeblue!5,
  colframe=pipeblue,
  fonttitle=\bfseries,
  title={Pipeline Step},
  #1
}

\newtcolorbox{codebox}[1][]{
  colback=codegray,
  colframe=gray!70,
  fonttitle=\bfseries\ttfamily,
  title={Code},
  #1
}

\newtcolorbox{termbox}[1][]{
  colback=termblack!5,
  colframe=termblack!60,
  fonttitle=\bfseries\ttfamily,
  title={Terminal},
  #1
}

% --- Research ---
\newtcolorbox{hypothesis}[1][]{
  colback=goldcolor!8,
  colframe=goldcolor,
  fonttitle=\bfseries,
  title={Hypothesis},
  #1
}

\newtcolorbox{resultbox}[1][]{
  colback=thmgreen!8,
  colframe=thmgreen,
  fonttitle=\bfseries,
  title={Expected Result},
  #1
}

% --- Constraints & Solutions ---
\newtcolorbox{constraintbox}[1][]{
  colback=constraintred!8,
  colframe=constraintred,
  fonttitle=\bfseries,
  title={Constraint},
  #1
}

\newtcolorbox{solutionbox}[1][]{
  colback=thmgreen!8,
  colframe=thmgreen,
  fonttitle=\bfseries,
  title={Solution},
  #1
}

\newtcolorbox{databox}[1][]{
  colback=pipeblue!5,
  colframe=pipeblue,
  fonttitle=\bfseries,
  title={Data Spec},
  #1
}

% --- Progress/Status ---
\newtcolorbox{successbox}[1][]{
  colback=lightgreen,
  colframe=thmgreen,
  fonttitle=\bfseries,
  title={Success},
  #1
}

\newtcolorbox{nextbox}[1][]{
  colback=lightorange,
  colframe=noteorange,
  fonttitle=\bfseries,
  title={Next Step},
  #1
}

% --- Fitting guide ---
\newtcolorbox{failbox}[1][]{
  colback=lightred,
  colframe=darkred,
  fonttitle=\bfseries,
  title={Failure Pattern},
  #1
}

\newtcolorbox{fixbox}[1][]{
  colback=lightgreen,
  colframe=thmgreen,
  fonttitle=\bfseries,
  title={Fix},
  #1
}

\newtcolorbox{lessonbox}[1][]{
  colback=lightorange,
  colframe=noteorange,
  fonttitle=\bfseries,
  title={Lesson Learned},
  #1
}

% --- Specification ---
\newtcolorbox{specbox}[1][]{
  colback=cvpurple!5,
  colframe=cvpurple,
  fonttitle=\bfseries,
  title={Specification},
  #1
}

\newtcolorbox{checklistbox}[1][]{
  colback=lightgreen,
  colframe=biogreen,
  fonttitle=\bfseries,
  title={Checklist},
  #1
}

% ========== Custom math commands ==========
% Vectors (lowercase)
\newcommand{\vx}{\mathbf{x}}
\newcommand{\vd}{\mathbf{d}}
\newcommand{\vo}{\mathbf{o}}
\newcommand{\vc}{\mathbf{c}}
\newcommand{\vr}{\mathbf{r}}
\newcommand{\vp}{\mathbf{p}}
\newcommand{\vv}{\mathbf{v}}
\newcommand{\vn}{\mathbf{n}}
\newcommand{\vu}{\mathbf{u}}
\newcommand{\vt}{\mathbf{t}}
\newcommand{\ve}{\mathbf{e}}
\newcommand{\vq}{\mathbf{q}}
% Vectors (uppercase)
\newcommand{\vF}{\mathbf{F}}
\newcommand{\vM}{\mathbf{M}}
\newcommand{\va}{\mathbf{a}}
\newcommand{\vg}{\mathbf{g}}
\newcommand{\vX}{\mathbf{X}}
% Bold Greek
\newcommand{\vmu}{\bm{\mu}}
\newcommand{\vbeta}{\bm{\beta}}
\newcommand{\vtheta}{\bm{\theta}}
\newcommand{\vpsi}{\bm{\psi}}
% Matrices
\newcommand{\mSigma}{\bm{\Sigma}}
\newcommand{\mR}{\mathbf{R}}
\newcommand{\mS}{\mathbf{S}}
\newcommand{\mW}{\mathbf{W}}
\newcommand{\mJ}{\mathbf{J}}
\newcommand{\mG}{\mathbf{G}}
\newcommand{\mT}{\mathbf{T}}
\newcommand{\mI}{\mathbf{I}}
\newcommand{\mB}{\mathbf{B}}
\newcommand{\mK}{\mathbf{K}}
\newcommand{\mP}{\mathbf{P}}
\newcommand{\mF}{\mathbf{F}}
\newcommand{\mE}{\mathbf{E}}
\newcommand{\mA}{\mathbf{A}}
\newcommand{\mH}{\mathbf{H}}
% Sets and groups
\newcommand{\RR}{\mathbb{R}}
\newcommand{\SO}{\mathrm{SO}}
% Units
\newcommand{\degps}{\,\text{deg/s}}
\newcommand{\Nm}{\ensuremath{\,\text{N}\cdot\text{m}}}
\newcommand{\BW}{\,\text{BW}}

% ========== Header/Footer ==========
\pagestyle{fancy}
\fancyhf{}
\fancyhead[L]{\leftmark}
\fancyhead[R]{\thepage}
\renewcommand{\headrulewidth}{0.4pt}

% ========== Title formatting ==========
\titleformat{\chapter}[display]
  {\normalfont\huge\bfseries\color{defblue}}
  {\chaptertitlename\ \thechapter}{20pt}{\Huge}

\titleformat{\section}
  {\normalfont\Large\bfseries\color{defblue!80!black}}
  {\thesection}{1em}{}

\titleformat{\subsection}
  {\normalfont\large\bfseries\color{defblue!60!black}}
  {\thesubsection}{1em}{}


\definecolor{solvable}{RGB}{30, 140, 60}
\definecolor{hard}{RGB}{200, 40, 40}
\definecolor{okgreen}{RGB}{20, 110, 60}
\definecolor{failred}{RGB}{200, 40, 40}

\begin{document}

\begin{titlepage}
\centering
\vspace*{1.5cm}
{\Large\color{gray} Experiment Log}
\vspace{0.5cm}

{\Huge\bfseries\color{defblue} 2026년 4월 12일\\[0.3cm]
실험 기록 및 결론}

\vspace{1cm}

{\LARGE\color{defblue!70!black} VGGT 캘리브레이션 $\rightarrow$ EasyMocap $\rightarrow$ SMPL 피팅\\
시도, 실패, 분석, 그리고 최종 결론}

\vspace{2cm}

\textbf{VGGT K,R,t 추출} $\xrightarrow{\checkmark}$
\textbf{삼각측량 999프레임} $\xrightarrow{\checkmark}$
\textbf{3D SMPL 피팅} $\xrightarrow{\times}$
\textbf{2D Reprojection으로 전환}

\vspace{2cm}
{\large 2026년 4월 12일}
\end{titlepage}

\tableofcontents
\newpage


% ====================================================================
\chapter{실험 1: VGGT 카메라 캘리브레이션}
% ====================================================================

\section{목적}
DA3의 K 역변환 문제 (비균일 스케일링, 가로/세로 혼재)를 우회하기 위해 VGGT로 카메라 파라미터를 새로 추정.

\section{환경}
\begin{itemize}[leftmargin=*]
  \item 서버: Elice Cloud, NVIDIA A100 80GB, Python 3.10, PyTorch 2.3.1+cu121
  \item VGGT: \texttt{facebook/VGGT-1B} (1B params, CVPR 2025 Best Paper)
  \item 입력: 7대 카메라 각 1프레임 (\texttt{00001.png}), cam1 portrait 90\textdegree{} 회전
  \item 모드: \texttt{crop} (pad 모드에서 $f_x/f_y$ 왜곡 발견 후 전환)
\end{itemize}

\section{결과}

\begin{table}[H]
\centering
\caption{VGGT 캘리브레이션 결과 (원본 해상도 기준)}
\renewcommand{\arraystretch}{1.3}
\begin{tabular}{l cc cc}
\toprule
\textbf{Camera} & $f_x$ & $f_y$ & $f_x/f_y$ & \textbf{해상도} \\
\midrule
cam1 (portrait) & 738.0 & 741.1 & 0.996 & 1080$\times$1920 \\
cam2 & 733.2 & 737.3 & 0.994 & 1920$\times$1080 \\
cam3 & 730.0 & 734.5 & 0.994 & 1920$\times$1080 \\
cam4 & 732.9 & 736.8 & 0.995 & 1920$\times$1080 \\
cam5 & 729.7 & 725.9 & 1.005 & 1920$\times$1080 \\
cam6 & 737.4 & 740.3 & 0.996 & 1920$\times$1080 \\
cam7 & 740.7 & 740.1 & 0.999 & 1920$\times$1080 \\
\midrule
\textbf{평균} & \textbf{734.6} & \textbf{736.6} & \textbf{0.997} & \\
\textbf{spread} & 11.0 & 14.4 & & \\
\bottomrule
\end{tabular}
\end{table}

\begin{summary}
$f_x/f_y$ 비율 0.994--1.005 (DA3 대비 4배 개선). $f_x$ spread 11px (동일 카메라 모델 일관성 확인). cx, cy가 정확히 이미지 중심. \textcolor{thmgreen}{\textbf{성공}}.
\end{summary}


% ====================================================================
\chapter{실험 2: EasyMocap 삼각측량}
% ====================================================================

\section{입력}
\begin{itemize}[leftmargin=*]
  \item \texttt{intri.yml}: VGGT K (7대)
  \item \texttt{extri.yml}: VGGT R, t (7대, Rodrigues + rotation matrix)
  \item \texttt{annots/\{1..7\}/}: OpenPose Body25 2D 키포인트 (999프레임)
\end{itemize}

\section{결과}
\begin{itemize}[leftmargin=*]
  \item 999프레임 삼각측량 완료 (1.4초)
  \item 25개 키포인트, confidence 0.72--0.93
  \item 3D 키포인트가 합리적인 인체 형태
\end{itemize}

\begin{table}[H]
\centering
\caption{삼각측량된 3D 키포인트 (Frame 0, 주요 관절)}
\renewcommand{\arraystretch}{1.2}
\begin{tabular}{l ccc c}
\toprule
\textbf{Joint} & $X$ & $Y$ & $Z$ & \textbf{Conf.} \\
\midrule
Nose & 0.133 & $-0.030$ & 0.404 & 0.93 \\
MidHip & $-0.053$ & 0.022 & 0.402 & 0.83 \\
RAnkle & $-0.252$ & 0.086 & 0.445 & 0.87 \\
LAnkle & $-0.272$ & 0.017 & 0.408 & 0.87 \\
\bottomrule
\end{tabular}
\end{table}

\begin{warningbox}
Nose-Ankle 거리 = 0.404 단위. 실제 피사체 키 183cm. \textbf{스케일 팩터 $\approx$ 4.53}. VGGT는 metric scale을 제공하지 않음.
\end{warningbox}

\begin{summary}
삼각측량 자체는 \textcolor{thmgreen}{\textbf{성공}}. 7대 카메라의 2D 키포인트가 VGGT K와 좌표계가 일치하여 정상 작동.
\end{summary}


% ====================================================================
\chapter{실험 3: EasyMocap SMPL 피팅 --- 1차 시도}
% ====================================================================

\section{설정}
\begin{itemize}[leftmargin=*]
  \item VGGT K, R, t (원본, 좌표계 수정 전)
  \item EasyMocap \texttt{triangulate\_fitSMPL.yml}: 삼각측량 $\rightarrow$ shape $\rightarrow$ RT $\rightarrow$ pose 순차 최적화
  \item SMPL\_NEUTRAL\_FIXED.pkl (21MB, full model)
\end{itemize}

\section{결과}

\begin{table}[H]
\centering
\caption{EasyMocap SMPL 피팅 결과 --- 1차 (좌표계 수정 전)}
\renewcommand{\arraystretch}{1.3}
\begin{tabular}{l cccc}
\toprule
\textbf{Frame} & $|\mR_h|$ (deg) & MPJPE & \textbf{Th} \\
\midrule
0 & 154.3 & 0.394 ($\sim$1456mm) & $[-0.20, 0.00, 0.40]$ \\
30 & 248.9 & 0.362 ($\sim$1339mm) & $[-0.32, -0.01, 0.49]$ \\
50 & 245.8 & 0.370 ($\sim$1369mm) & $[-0.32, -0.01, 0.49]$ \\
\bottomrule
\end{tabular}
\end{table}

\begin{warningbox}
$|\mR_h| = 248$\textdegree{} --- SMPL이 거의 뒤집어진 상태. MPJPE $\sim$1400mm --- 전혀 맞지 않음. \textcolor{hard}{\textbf{실패}}.
\end{warningbox}

\subsection{원인 분석}

삼각측량된 3D 키포인트에서 사람의 키 방향:
\begin{align*}
\text{Nose} - \text{MidHip} &= [0.186, -0.052, 0.002] \quad \rightarrow \text{주로 } +X \text{ 방향}
\end{align*}

SMPL의 키 방향: $+Y$ (Y-up).

\begin{keyidea}
VGGT 월드 좌표계에서 사람이 \textbf{X축 방향으로 서 있는데}, SMPL은 \textbf{Y축 방향}을 기대. EasyMocap이 이 불일치를 $\mR_h$로 보정하려 하면서 248\textdegree{} 회전이 발생하고, optimizer가 local minimum에 빠짐.
\end{keyidea}


% ====================================================================
\chapter{실험 4: 좌표계 수정 후 재시도}
% ====================================================================

\section{수정 내용}

삼각측량된 스켈레톤에서 사람의 방향 벡터를 추출하여 회전 행렬 $\mR_\text{fix}$ 계산:
\begin{align}
\hat{\mathbf{up}} &= \text{normalize}(\text{Nose} - \text{MidHip}) \\
\hat{\mathbf{right}} &= \text{normalize}(\text{RSho} - \text{LSho}) \\
\hat{\mathbf{fwd}} &= \hat{\mathbf{right}} \times \hat{\mathbf{up}} \\
\mR_\text{fix} &= [\hat{\mathbf{right}} \mid \hat{\mathbf{up}} \mid \hat{\mathbf{fwd}}]^\top
\end{align}

$\mR_\text{fix}$를 모든 카메라 extrinsics에 적용: $\mR_\text{cam}^\text{new} = \mR_\text{cam}^\text{old} \cdot \mR_\text{fix}^\top$

\section{검증}
수정 후 삼각측량된 키포인트:
\begin{align*}
\text{Nose} - \text{MidHip} &= [0.010, 0.192, -0.015] \quad \rightarrow \text{주로 } +Y \text{ (Y-up 확인)}
\end{align*}

\section{결과}

\begin{table}[H]
\centering
\caption{EasyMocap SMPL 피팅 결과 --- 좌표계 수정 후}
\renewcommand{\arraystretch}{1.3}
\begin{tabular}{l cccc}
\toprule
\textbf{Frame} & $|\mR_h|$ (deg) & MPJPE & 비고 \\
\midrule
0 & 143.2 & 0.374 ($\sim$1384mm) & 여전히 큰 회전 \\
30 & 175.6 & 0.447 ($\sim$1654mm) & 악화 \\
50 & 173.5 & 0.450 ($\sim$1665mm) & 악화 \\
\bottomrule
\end{tabular}
\end{table}

\begin{warningbox}
좌표계를 수정했음에도 $|\mR_h| \sim 143$--$175$\textdegree, MPJPE $\sim$1400mm. \textcolor{hard}{\textbf{여전히 실패}}.

원인: EasyMocap의 초기화가 VGGT의 정규화된 스케일 (사람 키 0.4 단위)에서 혼란에 빠짐. SMPL은 미터 단위 ($\sim$1.7m)인데 데이터는 0.4 단위이므로, optimizer가 스케일 차이를 $\mR_h$로 흡수하면서 잘못된 local minimum에 수렴.
\end{warningbox}


% ====================================================================
\chapter{실험 5: 직접 3D SMPL 피팅 (올바른 초기화)}
% ====================================================================

\section{설정}
EasyMocap을 우회하고 smplx + PyTorch로 직접 피팅:
\begin{itemize}[leftmargin=*]
  \item 스켈레톤에서 $\mR_h$ 초기값 계산 ($|\mR_h| = 5.5$\textdegree{} --- 정상)
  \item $\vt_h$ = MidHip 위치로 초기화
  \item 3-stage: RT(200 steps) $\rightarrow$ body pose(500 steps) $\rightarrow$ betas(300 steps)
  \item $\lambda_\theta = 0.001$, Adam optimizer
\end{itemize}

\section{결과}

\begin{table}[H]
\centering
\caption{직접 3D SMPL 피팅 결과 (Frame 0)}
\renewcommand{\arraystretch}{1.3}
\begin{tabular}{l cc}
\toprule
\textbf{Joint} & \textbf{Error} & \textbf{$\sim$mm (scale$\times$3700)} \\
\midrule
Nose & 0.096 & 355 \\
RElb & 0.012 & \textbf{44} \\
LWri & 0.024 & 88 \\
MidHip & 0.026 & 96 \\
LKne & 0.130 & 479 \\
LAnk & 0.135 & \textbf{499} \\
\midrule
\textbf{Mean (15 joints)} & \textbf{0.065} & \textbf{241} \\
\bottomrule
\end{tabular}
\end{table}

\begin{summary}
EasyMocap (1384mm) $\rightarrow$ 직접 피팅 (\textbf{241mm}). 5.7배 개선. 그러나 여전히 241mm는 부적합.
\begin{itemize}[leftmargin=*]
  \item 상체 (RElb 44mm, LWri 88mm) --- 비교적 양호
  \item 하체 (LKne 479mm, LAnk 499mm) --- 매우 나쁨
  \item 원인: \textbf{스케일 불일치}. SMPL 다리 길이 $\sim$0.8m vs 데이터 다리 범위 $\sim$0.2 단위
\end{itemize}
\end{summary}


% ====================================================================
\chapter{실험 6: 비등방 스케일 분석}
% ====================================================================

\section{발견}

SMPL T-pose와 삼각측량 데이터의 신체 비율 비교:

\begin{table}[H]
\centering
\caption{SMPL vs 삼각측량 신체 비율}
\renewcommand{\arraystretch}{1.3}
\begin{tabular}{l cc c}
\toprule
\textbf{측정} & \textbf{SMPL (m)} & \textbf{삼각측량} & \textbf{비율} \\
\midrule
키 (Nose-Ankle) & 0.090 & 0.404 & 0.22 \\
어깨 너비 & 0.646 & 0.099 & 6.52 \\
엉덩이 너비 & 0.566 & 0.065 & 8.72 \\
\bottomrule
\end{tabular}
\end{table}

\begin{warningbox}
비율이 0.22 $\sim$ 8.72로 \textbf{30배 이상 차이}. 이건 단순 스케일 문제가 아니라 \textbf{비등방 왜곡}:
\begin{itemize}[leftmargin=*]
  \item 키 방향 (Y축): 데이터가 SMPL보다 4.5배 큼
  \item 좌우 방향 (X축): SMPL이 데이터보다 6.5배 큼
\end{itemize}
\end{warningbox}

\begin{keyidea}
삼각측량 결과는 정확하지만, VGGT의 정규화된 스케일에서 \textbf{비율이 보존되지 않음}. 이는 3D 공간에서 SMPL을 맞추는 모든 방법의 근본적 한계.
\end{keyidea}


% ====================================================================
\chapter{최종 결론: 왜 2D Reprojection이 답인가}
% ====================================================================

\section{3D 경유 방식의 실패 패턴}

\begin{table}[H]
\centering
\caption{모든 3D 기반 SMPL 피팅 시도 결과}
\renewcommand{\arraystretch}{1.4}
\small
\begin{tabular}{c p{4.5cm} c p{4.5cm}}
\toprule
\textbf{\#} & \textbf{시도} & \textbf{MPJPE} & \textbf{실패 원인} \\
\midrule
1 & DA3 K + EasyMocap & --- & K 좌표계 불일치, Y/Z 스왑 \\
2 & 이전 직접 3D 피팅 & 219mm & 경량 SMPL, 부족한 iteration \\
3 & VGGT K + EasyMocap & 1384mm & 스케일 + 초기화 실패 \\
4 & VGGT K + 좌표계 수정 & 1384mm & 여전히 local minimum \\
5 & VGGT K + 직접 피팅 & 241mm & 비등방 스케일 왜곡 \\
\bottomrule
\end{tabular}
\end{table}

\begin{warningbox}[title={3D 경유 방식의 구조적 한계}]
2D $\rightarrow$ \textbf{삼각측량} $\rightarrow$ 3D $\rightarrow$ SMPL 피팅 경로에서 발생하는 문제:
\begin{enumerate}[leftmargin=*]
  \item \textbf{스케일}: VGGT는 metric scale 미제공. 삼각측량 결과가 임의 단위.
  \item \textbf{좌표계}: VGGT 월드 좌표 $\neq$ SMPL 좌표 규약. 변환 필요.
  \item \textbf{비등방 왜곡}: 축별로 스케일이 다르게 왜곡될 수 있음.
  \item \textbf{부족결정계}: 15관절 $\times$ 3좌표 = 45 방정식 $<$ 85 미지수.
  \item \textbf{초기화 민감}: 위 문제들로 인해 optimizer가 local minimum에 빠짐.
\end{enumerate}
\end{warningbox}

\section{2D Reprojection 방식의 구조적 이점}

\begin{keyidea}[title={2D Reprojection SMPL 피팅}]
\begin{equation}
\min_{\vtheta, \vbeta, \mR_h, \vt_h} \sum_{i=1}^{7} \sum_{j=1}^{25} w_j^i \left\| \pi(\mK_i, \mR_i, \vt_i, J_j(\vtheta, \vbeta, \mR_h, \vt_h)) - \vu_j^i \right\|^2
\end{equation}
여기서:
\begin{itemize}[leftmargin=*]
  \item $\pi$: 3D $\rightarrow$ 2D 투영 함수 (K, R, t 사용)
  \item $J_j$: SMPL forward kinematics로 계산된 3D 관절
  \item $\vu_j^i$: 카메라 $i$에서 OpenPose가 검출한 2D 키포인트
  \item $w_j^i$: OpenPose confidence
\end{itemize}
\end{keyidea}

\begin{table}[H]
\centering
\caption{3D 경유 vs 2D Reprojection 비교}
\renewcommand{\arraystretch}{1.4}
\begin{tabular}{p{4cm} p{5cm} p{5cm}}
\toprule
& \textbf{3D 경유 (실패)} & \textbf{2D Reprojection (다음 단계)} \\
\midrule
경로 & 2D $\rightarrow$ 삼각측량 $\rightarrow$ 3D $\rightarrow$ SMPL & 2D $\rightarrow$ SMPL (직접) \\
\midrule
방정식 수 & 15 $\times$ 3 = \textbf{45} & 25 $\times$ 7 $\times$ 2 = \textbf{350} \\
미지수 & 85 & 85 \\
결정성 & \textcolor{hard}{부족결정} (0.53$\times$) & \textcolor{solvable}{\textbf{과결정}} (4.1$\times$) \\
\midrule
스케일 문제 & \textcolor{hard}{있음} (metric 필요) & \textcolor{solvable}{\textbf{없음}} (픽셀 단위) \\
좌표계 문제 & \textcolor{hard}{있음} (Y-up/down) & \textcolor{solvable}{\textbf{없음}} (투영이 처리) \\
비등방 왜곡 & \textcolor{hard}{있음} & \textcolor{solvable}{\textbf{없음}} \\
\midrule
삼각측량 필요 & Yes & \textbf{No} \\
초기화 민감도 & 높음 & 낮음 (과결정) \\
\bottomrule
\end{tabular}
\end{table}

\begin{summary}[title={핵심 결론}]
\textbf{OpenPose 3D 삼각측량이 잘 되는 이유}와 \textbf{같은 이유}로 2D reprojection SMPL 피팅이 잘 될 것이다:
\begin{itemize}[leftmargin=*]
  \item 삼각측량: 7대 카메라의 2D 정보를 \textbf{직접} 사용 $\rightarrow$ 과결정, 선형, 안정적
  \item 2D reproj: 7대 카메라의 2D 정보를 \textbf{직접} 사용 $\rightarrow$ 과결정, 비선형이지만 안정적
  \item 3D 경유: 2D $\rightarrow$ 3D 변환 과정에서 스케일/좌표계 정보 손실 $\rightarrow$ 부족결정, 불안정
\end{itemize}

\textbf{다음 단계}: 7대 카메라의 OpenPose 2D 키포인트 + VGGT K, R, t를 사용하여 2D reprojection loss로 SMPL을 직접 피팅하는 스크립트를 구현하고 실행한다.
\end{summary}


\section{추가 확인 사항}

\begin{itemize}[leftmargin=*]
  \item \textbf{피사체 키}: 183cm (확인됨) --- 2D 방식에서는 불필요하지만 검증용으로 사용
  \item \textbf{체커보드}: 없음 --- focal length GT 불가, VGGT K의 정확도는 reprojection error로 간접 검증
  \item \textbf{SMPL 모델}: SMPL\_NEUTRAL\_FIXED.pkl (21MB, full) --- 서버에서 이미 EasyMocap에 연결됨
  \item \textbf{GART 다시점 확장}: \texttt{training\_view} 1줄 변경으로 가능 (코드 분석 완료)
  \item \textbf{GART $\rightarrow$ OpenSim}: SMPL2AddBiomechanics로 연결 가능 (오픈소스 존재)
\end{itemize}


\end{document}
